\begin{abstract}
Arbitrary-scale super-resolution (ASISR) predicts pixels at continuous coordinates, but recovery above the low-resolution Nyquist frequency and the role of scale-dependent readout remain difficult to assess. We study frequency recovery and signed paired readout losses, separating input sampling, target measurement and network cell conditioning. Across three datasets, normalising frequency by the sampling scale reduces variation in all nine evaluated model--dataset groups; the pre-specified collapse tolerance nevertheless passes only six. Evaluated decoders exhibit similar recovery in 10 of 12 low-reference-recovery cells, without establishing a universal decoder limit. Two further training seeds support coarse readout-loss ordering in all four architecture--seed groups. The three-valued footprint predictor has the same Spearman ranks as band index, so this result does not validate its quantitative form. Averaging a full model over training-scale cell inputs predicts loss ordering, but fails function similarity in all four groups and improves on the footprint score in only two. Additional frozen-weight probes improve integral-forecast magnitude for SC-INR but miss negative benefits; neither tested forecast meets the complete criterion. Negative readout benefits and these failed stronger tests delimit what the baseline establishes: empirical scale and frequency structure, rather than a shared integral-readout mechanism or an adaptive model-selection rule.
\end{abstract}
